\documentclass[11pt]{article}

\usepackage[margin=1in]{geometry}
\usepackage{amsmath,amssymb,amsthm}
\usepackage{enumitem}
\usepackage{graphicx}
\usepackage[hidelinks]{hyperref}

\newtheorem{theorem}{Theorem}
\newtheorem{lemma}{Lemma}
\newcommand{\J}{\mathcal J}
\newcommand{\Fin}{\mathrm{Fin}}
\newcommand{\limsupJ}{\J\text{-}\limsup}
\newcommand{\limJ}{\J\text{-}\lim}

\title{A Rapid+ P+ Extension of Leonetti's Sliding-Jump Theorem}
\author{Run \texttt{fa\_banach\_001}, target arXiv:2201.13059}
\date{2026-06-30}

\begin{document}
\maketitle

\section*{Source Problem}

The source paper is Paolo Leonetti's arXiv:2201.13059,
\emph{Regular matrices of unbounded linear operators}. In the closing remarks,
page 24, the source asks whether Theorem 4.1 holds for the larger class of
rapid+ P+-ideals, whether this condition characterizes that class, and whether
an analogous question holds for Theorem 3.6.

\begin{center}
\includegraphics[width=0.98\linewidth]{figures/open_problem_crop.png}
\end{center}

\section*{Definitions and Source Theorem}

For an ideal $\J$ on $\omega$, write $\J^+=\mathcal P(\omega)\setminus\J$ and
$\J^\star=\{\omega\setminus A:\ A\in\J\}$.

The ideal $\J$ is P+ if every decreasing sequence $(S_n)$ in $\J^+$ has a
pseudointersection $S\in\J^+$ such that $S\setminus S_n$ is finite for every
$n$. It is rapid+ if, for every $S\in\J^+$ and every infinite
$F\subset\omega$, there is $S'\subset S$ with $S'\in\J^+$ and
\[
  |S'\cap[0,m]|\le |F\cap[0,m]| \quad (m\in\omega).
\]

Theorem 4.1 of the source says the following when $\J$ is countably generated.
Let $A=(A_{n,k})$ be a matrix of linear operators in $\mathcal L(X,Y)$ satisfying
the source conditions (T1$^\flat$), (T3$^\natural$), and (T6$^\flat$). Then
there is a sequence $\mathbf x=(x_k)$ with values in the unit sphere $S_X$ such
that
\[
  \limsupJ_n \|A_{n,\omega}\|
  =
  \limsupJ_n \|A_n\mathbf x\|.
\]
Here $A_{n,E}$ denotes the finite-or-infinite block operator used in the source,
and $A_{n,\omega}$ denotes the full row operator.

\section*{Idea of the Proof}

Leonetti's proof already contains almost all of the argument. It constructs
decreasing $\J$-positive ``good'' sets $S_n$ and then needs a $\J$-positive
sequence of row indices $s_n$ with $s_n\in S_n$. Countable generation supplies
such a sequence by making the chosen set escape every generator.

For rapid+ P+ ideals there is a replacement selection principle. P+ gives a
positive set $S$ almost contained in every $S_n$. For each $n$, only finitely
many points of $S$ fail to be in $S_n$. Choose an infinite comparison set $F$
whose $n$th element lies beyond the first $n$ finite exceptional sets. Rapid+
then thins $S$ to a positive subset whose increasing enumeration is at least as
sparse as $F$. Its $n$th element is therefore in $S_n$. With this selected
sequence, the original sliding-jump estimates are unchanged.

\begin{lemma}[rapid+ P+ diagonal thinning]\label{lem:thin}
Let $\J$ be a rapid+ P+ ideal and let $(S_n)$ be a decreasing sequence in
$\J^+$. Then there are a set $R=\{r_0<r_1<\cdots\}\in\J^+$ and an integer
$n_0$ such that $r_n\in S_n$ for every $n\ge n_0$.
\end{lemma}

\begin{proof}
By P+, choose $S\in\J^+$ such that $S\setminus S_n$ is finite for every $n$.
For each $n$, let
\[
  b_n=1+\max\bigcup_{i\le n}(S\setminus S_i),
\]
with the convention that the maximum of the empty set is $0$. Increasing
$b_n$ if necessary, choose a strictly increasing infinite set
$F=\{f_0<f_1<\cdots\}$ with $f_n\ge b_n$ for every $n$.

By rapid+, there is $R\subset S$ with $R\in\J^+$ and
$|R\cap[0,m]|\le |F\cap[0,m]|$ for all $m$. Let
$R=\{r_0<r_1<\cdots\}$. This counting inequality implies $r_n\ge f_n$ for every
$n$: otherwise, if $r_n<f_n$, then at $m=r_n$ the set $R$ has at least $n+1$
points while $F$ has at most $n$ points. Thus $r_n\ge f_n\ge b_n$. Since
$r_n\in S$ and $r_n$ is larger than every point of $S\setminus S_n$, we have
$r_n\in S_n$ for all $n$.
\end{proof}

\begin{theorem}
The conclusion of Leonetti's Theorem 4.1 remains true if the hypothesis
``$\J$ is countably generated'' is replaced by ``$\J$ is rapid+ and P+''.
\end{theorem}

\begin{proof}
We follow the notation of the source proof. Let
\[
  \eta_0=\limsupJ_n\|A_{n,\omega}\|.
\]
As in the source proof, (T6$^\flat$) implies that finite initial blocks are
$\J$-null, and together with (T1$^\flat$) this gives $\eta_0<\infty$. If
$\eta_0=0$, any sequence in $S_X^\omega$ works. Assume $\eta_0>0$.

The source proof defines decreasing sets
\[
 E_n=\left\{t\in J_2:\left|\|A_{t,\omega}\|-\eta_0\right|
          \le \frac{\eta_0}{2^n}\right\}\in\J^+
\]
and, after intersecting with finite-block control sets from (T6$^\flat$),
obtains a decreasing sequence of good sets $S_n\in\J^+$. In the countably
generated proof, one chooses $s_n\in S_n$ escaping the generators. Replace that
choice by Lemma \ref{lem:thin}: pass to $R=\{r_0<r_1<\cdots\}\in\J^+$ with
$r_n\in S_n$ for all sufficiently large $n$, and discard the finite initial
segment. Relabel the tail as $(s_n)$.

Now choose the increasing integers $m_n$ exactly as in the source proof, using
(T3$^\natural$) to ensure
\[
  \|A_{s_n,\ge m_n}\|\le \frac{\eta_0}{2^n},
\]
and choose vectors $x_k\in S_X$ on the blocks
$M_n=(m_{n-1},m_n]$ so that
\[
  \left\|\sum_{k\in M_n}A_{s_n,k}x_k\right\|
  \ge \|A_{s_n,M_n}\|-\frac{\eta_0}{2^n}.
\]
Because $s_n\in S_n$, the same estimates as in the source give, for all large
$n$,
\[
  \|A_{s_n}\mathbf x\|
  \ge
  \eta_0\left(1-\frac1{2^{\,n-3}}\right).
\]
The selected set $\{s_n:\ n\in\omega\}$ is in $\J^+$, so
\[
  \limsupJ_n\|A_n\mathbf x\|\ge \eta_0.
\]
The reverse inequality is the elementary row domination inequality used at the
start of the source proof:
\[
  \|A_n\mathbf x\|\le \|A_{n,\omega}\| \quad (n\in\omega).
\]
Hence equality holds.
\end{proof}

\section*{Verification Notes}

Only the selection step in the proof of source Theorem 4.1 changes. All analytic
estimates, including the finite-column null estimate, tail control from
(T3$^\natural$), and the block-vector choice, are verbatim from Leonetti's
proof. The finite initial segment lost in Lemma \ref{lem:thin} is irrelevant to
$\J$-limsup because finite sets belong to any ideal on $\omega$.

\section*{Limitations and Novelty Check}

This proves the first extension asked in the closing remarks. It does not prove
that Theorem 4.1 characterizes rapid+ P+-ideals, and it does not settle the
analogous question for Theorem 3.6 or the matrix class
$(c(X,\mathcal Z),c(Y,\mathcal Z))$.

The run indexes were searched for arXiv:2201.13059 and for terms including
``regular matrices'', ``unbounded linear operators'', ``rapid+ P+'', and
``sliding jump''. A bounded web search for the source title together with
rapid+ P+ and Theorem 4.1 found no exact later answer.

\section*{References}

\begin{enumerate}
\item P. Leonetti, \emph{Regular matrices of unbounded linear operators},
arXiv:2201.13059.
\item B. De Bondt and H. Vernaeve, \emph{Filter-dependent versions of the
uniform boundedness principle}, J. Math. Anal. Appl. 495 (2021), Paper
No. 124705.
\end{enumerate}

\end{document}
